% ECONOMICS_PROBLEM: {"id": "C31-129", "title": "Causal inference under interference and network spillovers", "jel_code": "C31", "source_id": "OP-0101"}
% !TeX program = xelatex
\documentclass[11pt,a4paper]{article}
\usepackage[margin=23mm]{geometry}
\usepackage{fontspec}
\setmainfont{Latin Modern Roman}
\usepackage{amsmath,amssymb,mathtools}
\usepackage{xcolor,enumitem,booktabs,longtable,array,xurl,hyperref}
\definecolor{muted}{HTML}{53636C}
\hypersetup{colorlinks=true,urlcolor=blue,linkcolor=blue}
\setlength{\parindent}{0pt}
\setlength{\parskip}{5pt}
\newcommand{\R}{\mathbb R}
\newcommand{\E}{\mathbb E}
\newcommand{\N}{\mathbb N}
\newcommand{\ind}{\mathbf 1}
\newcommand{\Var}{\operatorname{Var}}
\newcommand{\Cov}{\operatorname{Cov}}
\newcommand{\supp}{\operatorname{supp}}
\DeclareMathOperator*{\argmin}{arg\,min}
\DeclareMathOperator*{\argmax}{arg\,max}
\newcommand{\entrymeta}[3]{{\sffamily\small\color{muted}\textbf{\texttt{#1}}\quad|\quad #2\par #3\par}}
\newcommand{\Problemref}[1]{\texttt{#1}}
\newcommand{\JELref}[2]{\texttt{#2} (JEL \texttt{#1})}
\begin{document}
\section*{Causal inference under interference and network spillovers}
\textbf{Problem ID:} \texttt{C31-129}\quad \textbf{JEL:} \texttt{C31}\par
% BEGIN ECONOMICS STATEMENT
\label{problem:OP-0101}
\label{atlas:prob:E1}

\noindent\textbf{Original identifier:} E1 (atlas item 101). \textbf{Field:} Econometrics and causal inference. \textbf{Classification:} editorial research family with established restricted solutions; current open status not certified.

\subsection*{Definitions and assumptions}

For a finite population of $n$ units, let $A\in\{0,1\}^{n\times n}$ be a pretreatment network with zero diagonal, $Z\in\{0,1\}^n$ the randomized assignment, and $\pi(z)=\Pr(Z=z\mid A,X)$ its known design conditional on pretreatment covariates $X$. Let $Y_i(z)$ denote unit $i$'s potential outcome under the entire assignment. A specified exposure map $h_i(z,A)$ takes values in a finite set $\mathcal E$. Exact exposure sufficiency means $Y_i(z)=y_i(h_i(z,A))$ for every assignment in the design support. Define $p_i(e)=\Pr\{h_i(Z,A)=e\}$ and, for two exposures observable with positive probability for every unit, $\tau(e,e')=n^{-1}\sum_i[y_i(e)-y_i(e')]$. This is a finite-population exposure contrast, not automatically a direct effect holding neighbors' treatment fixed.

\subsection*{Formal research question}

Characterize design, network and exposure-class restrictions under which $\tau(e,e')$ is identified and admits honest inference when $h_i$ is learned. With a known sufficient map and $p_i(e)>0$, the Horvitz--Thompson estimator
\[
\widehat\mu(e)=\frac1n\sum_{i=1}^n
\frac{\mathbf1\{h_i(Z,A)=e\}Y_i(Z)}{p_i(e)}
\]
is design-unbiased for $n^{-1}\sum_i y_i(e)$. For a misspecified map, instead define $y_i(e)=\mathbb E_\pi[Y_i(Z)\mid h_i(Z,A)=e]$; this fixes a design-specific exposure contrast rather than a stable response shared by all assignments. The frontier is to obtain a confidence set $C_n$ satisfying $\liminf_{n\to\infty}\inf_{P\in\mathcal P_n}\Pr_P\{\tau_P\in C_n\}\ge1-\alpha$, for a stated $\alpha\in(0,1)$ and class $\mathcal P_n$ containing bounded outcomes, overlapping exposure probabilities, specified dependence growth, and bounded exposure misspecification. Establish which rates of network degree, probability decay and map complexity permit informative diameter; otherwise derive a lower bound or identified set.

\subsection*{Interpretation and scope}

A learned map changes both the estimand and its assignment probabilities. Independent map training, design-valid conditioning or simultaneous inference over maps must account for this selection. If network formation changes with treatment, replace $A$ by potential networks $A(z)$ and specify an intervention on formation; conditioning on the realized network can create selection bias. Without exposure restrictions, one realized assignment cannot recover all $2^n$ potential outcomes. Sharp-null randomization tests remain different from confidence intervals for heterogeneous average effects.

\subsection*{Source audit and status}

[S1] concerns the reflection problem for simultaneous social outcomes; it is related background, not the source of an arbitrary-interference theorem. [S2] treats partial interference; [S3] provides established design-based methods using exposure mappings. Accordingly the unrestricted original question is a research family with substantial solved subcases. The formal extension above is editorial and the citations do not certify that this exact extension remains unresolved in October 2026.

\subsection*{References}

\begin{enumerate}[label={[\textnormal{S}\arabic*]},leftmargin=*]

\item Charles F. Manski (1993). \emph{Identification of Endogenous Social Effects: The Reflection Problem}. Review of Economic Studies 60(3), 531–542. \url{https://doi.org/10.2307/2298123}. \textbf{Locator:} Publisher or institutional abstract and bibliographic record. \textbf{Support and limits:} Simultaneous social effects identification; not a theorem about arbitrary treatment interference.

\item Michael G. Hudgens and M. Elizabeth Halloran (2008). \emph{Toward Causal Inference With Interference}. Journal of the American Statistical Association 103(482), 832–842. \url{https://doi.org/10.1198/016214508000000292}. \textbf{Locator:} Publisher or institutional abstract and bibliographic record. \textbf{Support and limits:} Defines direct, indirect, total and overall effects with partial interference; does not identify all endogenous-network effects.

\item Peter M. Aronow and Cyrus Samii (2017). \emph{Estimating Average Causal Effects Under General Interference, with Application to a Social Network Experiment}. Annals of Applied Statistics 11(4), 1912–1947. \url{https://doi.org/10.1214/16-AOAS1005}. \textbf{Locator:} Publisher or institutional abstract and bibliographic record. \textbf{Support and limits:} Design-based exposure-mapping estimation under known assignment probabilities; not unrestricted unknown interference.

\end{enumerate}

\subsection*{Common conventions: Source mathematical conventions}

Unless an entry states otherwise, agent and firm sets are finite. State and action spaces are measurable, strategies and outcome maps are measurable, probability laws are specified, and expectations exist for the targets under discussion. Infinite-horizon discount factors lie in $(0,1)$ and discounted objectives are finite. Norms, tolerances and complexity input sizes must be specified for computational comparisons. Constants or primitives that appear locally are inputs, not empirical facts asserted by this catalogue.

\textbf{Identification.} A statistical model $m\in\mathcal M$ implies an observable law $P_m$ and target $\theta(m)$. At law $P$, its identified set is
\[
\Theta(P;\mathcal M)=\{\theta(m):m\in\mathcal M,\ P_m=P\}.
\]
Point identification holds when this set is a singleton, or equivalently when $P_{m_1}=P_{m_2}$ implies $\theta(m_1)=\theta(m_2)$ for all compatible models. Sharp bounds mean bounds attained, or approached when the boundary is not attained, by compatible models. Writing this set is a definition, not a solution: the substantive task is to establish informative restrictions and derive or estimate the set. An unrestricted-confounding model may give only trivial bounds.

\textbf{Inference.} Identification, estimability and valid uncertainty quantification are separate questions. Fix $\alpha\in(0,1)$. A confidence procedure $C_n$ over a declared law class $\mathcal P$ has asymptotic uniform coverage at level $1-\alpha$ when
\[
\liminf_{n\to\infty}\inf_{P\in\mathcal P}\Pr_P\{\theta(P)\in C_n\}\geq1-\alpha.
\]
For set-identified targets the coverage event must be defined separately, for example coverage of the full identified set. If an entry uses identification-indexed models instead of uniquely indexed laws, the same coverage convention is applied over those models. Pointwise limits do not establish uniform validity, and asymptotic validity does not guarantee finite-sample precision. Sample sizes, dependence structures and weak-identification sequences are part of each application.

\textbf{Counterfactuals.} $Y_i(a)$ is a potential outcome under a specified intervention, including its timing, implementation and versions. It is not $Y_i$ conditional on voluntarily choosing $a$. A full intervention vector permits interference; shorthand scalar treatment notation does not silently impose no interference. Assignment, selection, attrition, anticipation and measurement must be specified. A claim about an instrument's local effect is not automatically a population policy effect.

\textbf{Economic equilibrium.} A counterfactual model includes best responses, constraints, feasibility, market clearing, state transitions and expectations consistent with the stated information. When several equilibria exist, either a declared selector is an input or the target is set-valued. Comparisons across policy regimes must use the stated selection convention. Reduced-form relationships are not presumed invariant after an equilibrium-changing intervention.

\textbf{Optimization and welfare.} Optimization is over the stated feasible set. If attainment is not established, $\sup$ or $\inf$ and $\varepsilon$-optimal choices replace an asserted maximizer or minimizer. For example, with value $V=\sup_{a\in\mathcal A}W(a)<\infty$, an $\varepsilon$-optimal set is $\{a\in\mathcal A:W(a)\geq V-\varepsilon\}$ for $\varepsilon>0$. Benefits, costs, risk and health or environmental outcomes can be aggregated only using declared units and conversion weights. Interpersonal utility weights, the treatment of fiscal transfers and foreign profits, discounting, and the behavioral welfare criterion are explicit inputs. A comparative-static derivative additionally requires existence and appropriate differentiability of the selected counterfactual.

\textbf{Sources.} Local labels [S1], [S2], and so forth restart in every question. Locators identify the relevant abstract, model, section or result at the level actually checked; a bibliographic or abstract-level verification is not represented as inspection of an entire proof. Working papers are distinguished from later publications and changing versions. Source summaries are paraphrased. All URL checks and editorial formulations refer to the audit date stated above.
% END ECONOMICS STATEMENT
\end{document}
